{
  \renewcommand{\sectionsubtitle}{Content distribution without a human operator}
  \section{Research Angle}
}


\begin{frame}{Centralized distribution means a single point of control}
  Mainstream content distribution concentrates control on one entity
  that owns both the hardware and the content.
  \begin{itemize}
    \item An operator decides what stays online. Availability depends
      on profitability, policies, and legality.
    \item Peer-to-peer protocols distribute the content
      across many nodes, but the hardware is
      still \emph{managed} by humans.
  \end{itemize}

  \begin{alertblock}{The open research gap}
    Can the \textbf{infrastructure layer itself} be operator-independent, with no human controlling it after genesis?
  \end{alertblock}
\end{frame}

%TODO: add here slide that very briefly details the high level goal of the system, how it fits in the open research gap and how it fixes the two drawbacks mentioned on the slide above. only then it leads into the research question below. 

{
  \setlength{\splitpos}{0.55\paperwidth} % colored pane left, title shifts right
  \leftfooterwhitetrue
  \begin{frame}{Research question}
    \begin{columns}[T]
      \begin{textcolumn}
        \color{white}\large
        Can a network of autonomous software agents, controlled only by inherited economic
        parameters, \textbf{survive} and \textbf{grow} without further human
        intervention?
      \end{textcolumn}
      \begin{textcolumn}
        The hypothesis under test: two outcomes must hold, network
        \textcolor{tud Orange}{persistence} and \textcolor{tud Orange}{expansion},
        with automated control and no human in the loop after genesis.
        
        Falsifiable: survival and growth, can be disproved by data.
      \end{textcolumn}
    \end{columns}
  \end{frame}
}
\setlength{\splitpos}{0pt}\leftfooterwhitefalse % reset stateful split/footer


{
\setbeamerfont{itemize/enumerate body}{size=\small} % template pins lists to \normalsize
\setbeamerfont{footnote}{size=\scriptsize} % shrink full-reference citation footnotes
\footerinfootnotestrue % full-reference footnotes -> move footer into footnote band
\begin{frame}{Where this sits scientifically}
  \small
  \begin{columns}[T, onlytextwidth]
    \begin{column}{0.48\textwidth}
      Prior self-replicators ran on simulated resources:
      \begin{itemize}
        \item von Neumann's automata\cite{vonneumann} and Ray's
          Tierra\cite{tierra}. Running on simulated CPU cycles
          and RAM.
        \item DAOs and autonomic computing self-manage, but only over
          virtual, quota-based resources.
      \end{itemize}

      \textcolor{tud Orange}{\textbf{Research Novelty:}} first self-replicating system on
      \textbf{real infrastructure paid for with real money}.
    \end{column}
    \begin{column}{0.48\textwidth}
      We propose three contributions:
      \begin{enumerate}
        \item \textbf{Engineering} --- first self-replicating economic agent that
          buys real infrastructure on an open market.
        \item \textbf{Method} --- unmodified production code run against a replica
          Bitcoin network and a mock VPS marketplace.
        \item \textbf{Empirical} --- fleet dynamics and one heritable trait under
          selection.
      \end{enumerate}
    \end{column}
  \end{columns}
\end{frame}
}
